\documentclass[conference]{IEEEtran}
\IEEEoverridecommandlockouts

\usepackage{cite}
\usepackage{amsmath,amssymb,amsfonts}
\usepackage{algorithmic}
\usepackage{graphicx}
\usepackage{textcomp}
\usepackage{xcolor}
\usepackage{hyperref}
\usepackage{booktabs}
\usepackage{array}
\usepackage{multirow}
\usepackage{microtype}

\graphicspath{{figures/}}

\hypersetup{
  colorlinks=true,
  linkcolor=black,
  citecolor=black,
  urlcolor=blue
}

\def\BibTeX{{\rm B\kern-.05em{\sc i\kern-.025em b}\kern-.08em
T\kern-.1667em\lower.7ex\hbox{E}\kern-.125emX}}

\begin{document}

\title{In-Context Tabular Regression for Training-Free Gene and Pathway  Prediction}

\author{
\IEEEauthorblockN{Ushashi Bhattacharjee\textsuperscript{1},
Alloy Das\textsuperscript{1},
Saria Hannan\textsuperscript{1},
Tirtho Roy\textsuperscript{1},
and Soumik Sarkar\textsuperscript{1}}
\IEEEauthorblockA{
\textsuperscript{1}Department of Mechanical Engineering,
Iowa State University, Ames, IA, USA\\
\{ushashi, alloydas, shannan, tirtho, soumiks\}@iastate.edu
}

}

\maketitle

% ============================================================================
\begin{abstract}
% ============================================================================
Predicting molecular expression from histopathology images is a fundamental
challenge in computational pathology, with direct applications to cancer
diagnosis and therapy design.
The PEaRL framework~\cite{pearl2025} advances this problem through
pathway-enhanced contrastive pretraining followed by supervised MLP heads
that regress gene and pathway expression from frozen image embeddings.
In this work we ask: can a pretrained in-context tabular foundation model ---
\textsc{TabPFN}~\cite{tabpfn2023} --- replace PEaRL's gradient-trained MLP
heads without any fine-tuning, while preserving or improving predictive
performance?
We evaluate on two HEST-1k cohorts~\cite{hest1k}: Breast cancer
(36 sections, 13{,}620 spots, 775 pathways) and Skin cancer
(12 sections, 8{,}671 spots, 609 pathways).
Under matched experimental conditions, TabPFN achieves gene PCC of $0.5203$
vs.\ $0.5196$ for the reproduced MLP on Breast (+0.14\%), and $0.8704$
vs.\ $0.8691$ on Skin (+0.15\%), while simultaneously reducing gene MSE by
8.0\% and 0.9\% respectively --- with \emph{zero gradient updates}.
Per-dimension PCC analysis further characterizes the qualitative differences
between the two heads.
All code, configurations, and SLURM workflows are released at
\url{https://github.com/tirtho149/PeARL-TabFPN.git}.
\end{abstract}

\begin{IEEEkeywords}
spatial transcriptomics, computational pathology, foundation models,
TabPFN, in-context learning, pathway analysis, HEST-1k, UNI,
contrastive learning, gene expression prediction
\end{IEEEkeywords}

% ============================================================================
\section{Introduction}
\label{sec:intro}
% ============================================================================

Spatial transcriptomics (ST) technologies such as 10x Visium simultaneously
capture gene expression and spatial position across tissue sections, enabling
molecular mapping of tissue architecture at near-cellular resolution~\cite{hest1k}.
However, ST assays are expensive, slow, and not routinely performed in clinical
settings.
Computational prediction of molecular profiles from standard H\&E histopathology
images would enable retrospective analysis of archived tissue banks, accelerate
drug-target discovery, and support cost-effective precision oncology.

Cancer diagnosis relies on tissue morphology, but disease progression is driven
by molecular function.
Bridging this gap --- linking what tissue \textit{looks like} to what it is
\textit{doing} at the molecular level --- requires models that capture
coordinated molecular programs rather than individual gene signals.
Pathways, defined as networks of molecular interactions where genes function
together, aggregate noisy per-gene signals into compact and biologically
coherent modules that align with histological tissue patterns.

The PEaRL framework~\cite{pearl2025} addresses this by replacing individual
gene regression targets with single-sample gene-set enrichment (ssGSEA)
pathway scores during contrastive pretraining, producing image embeddings that
are semantically aligned with biological process structure.
On HEST-1k Breast, PEaRL achieves gene PCC of $0.587 \pm 0.036$ and pathway
PCC of $0.506 \pm 0.027$, outperforming prior methods including
ST-Net~\cite{stnet}, His2ST~\cite{his2st}, BLEEP~\cite{bleep}, and
MCLSTExp~\cite{mclstexp}.

A complementary line of work has shown that \textsc{TabPFN}~\cite{tabpfn2023},
a transformer meta-trained on millions of synthetic tabular tasks, matches or
exceeds gradient-boosted trees on small-to-medium tabular datasets with a
single forward pass and no per-task hyperparameter tuning.
After PEaRL's Stage-1 contrastive pretraining, the image embeddings form a
$(N \times 256)$ structured tabular matrix --- precisely the setting where
TabPFN's Bayesian inductive prior is most effective.

This raises a concrete question:
\textit{does TabPFN's in-context prior transfer to the structured biological
embeddings produced by PEaRL's contrastive pretraining?}
Our contributions are:
\begin{itemize}
    \item A \textbf{faithful PEaRL reproduction} on HEST-1k Breast and Skin
          cohorts with paper-matched preprocessing (8-neighbor spatial smoothing,
          Reactome + MSigDB Hallmark pathway pool, ssGSEA, Scanpy-flavor HVG
          selection).
    \item A \textbf{drop-in head replacement} in \textit{pure} mode: one
          \texttt{TabPFNRegressor} (32 estimators) per output dimension, covering
          all 1{,}775 (Breast) or 1{,}609 (Skin) dimensions with zero
          gradient updates.
    \item \textbf{Multi-dataset quantitative evaluation} on Breast and Skin
          cohorts reporting PCC, MSE, and MAE for both gene and pathway targets.
    \item \textbf{Qualitative analysis} of per-dimension PCC distributions
          across all gene and pathway output dimensions.
    \item A fully \textbf{reproducible open-source release} with SLURM workflows
          for Iowa State's Nova HPC cluster.
\end{itemize}

% ============================================================================
\section{Related Work}
\label{sec:related}
% ============================================================================

\subsection{Spatial Gene Expression Prediction}

ST-Net~\cite{stnet} first demonstrated CNN-based prediction of gene expression
from Visium breast cancer patches.
His2ST~\cite{his2st} combined Vision Transformers with graph neural networks
to aggregate neighboring spot information.
BLEEP~\cite{bleep} and MCLSTExp~\cite{mclstexp} incorporated contrastive
learning between image and expression embeddings.
A common limitation of all these methods is that they regress raw per-gene
values (typically 15{,}000+ dimensions), which are sparse and technically noisy;
using only a small subset of highly variable genes (HVGs) further limits
predictive scope.

PEaRL~\cite{pearl2025} overcomes this by using ssGSEA pathway scores ---
which aggregate coordinated gene signals into compact modules --- as the
contrastive pretraining target, producing representations that are semantically
coherent with biological processes visible in histological images.

\subsection{Foundation Models for Pathology}

UNI~\cite{uni2024} is a ViT-Large model trained via DINOv2 on 100{,}000 H\&E
WSIs, producing 1024-d CLS embeddings that serve as the strongest
publicly-available spot-level encoder on HEST-1k.
CONCH~\cite{conch2024} extends this with a vision-language formulation.
The HEST-1k benchmark~\cite{hest1k} consolidated multi-organ Visium data with
aligned WSIs, enabling systematic cross-method comparison.

\subsection{Pretrained Tabular Models}

TabPFN~\cite{tabpfn2023} is a transformer meta-trained on synthetically-generated
tabular tasks. It performs in-context prediction: a single forward pass on
concatenated (training-set, query-point) pairs returns a predictive distribution
without gradient updates.
The v2 release incorporates automatic feature interaction modeling and ensemble
averaging, improving performance on small and medium tabular problems.
To our knowledge, no prior work has applied a pretrained tabular foundation
model as a drop-in prediction head for spatial transcriptomics.

% ============================================================================
\section{Methodology}
\label{sec:method}
% ============================================================================

\subsection{PEaRL Framework Overview}

Figure~\ref{fig:pearl_arch} illustrates the original PEaRL architecture.
Each Visium spot contributes a $224\!\times\!224$ RGB histology patch, a
gene expression vector (CP10K + log-$(1{+}x)$ + 8-neighbor smoothing),
and a 2D spatial coordinate.
Training proceeds in two sequential stages.

\begin{figure}[t]
\centering
\includegraphics[width=\columnwidth]{Screenshot 2026-05-14 at 11.03.43 PM.png}
\caption{PEaRL framework~\cite{pearl2025}.
\textbf{Training}: image patches pass through UNI (Image Encoder) to produce
$H_\text{image} \in \mathbb{R}^{256}$; ssGSEA pathway scores + spatial coords
pass through the Pathway Encoder to produce $H_\text{path} \in \mathbb{R}^{256}$.
Both are aligned via NT-Xent contrastive loss $\mathcal{L}_{CL}$.
\textbf{Inference}: the trained Image Encoder feeds two MLP heads (Gene MLP,
Pathway MLP) for downstream expression prediction.}
\label{fig:pearl_arch}
\end{figure}

\textbf{Stage~1 (Contrastive Pretraining).}
A pathway encoder maps each spot's ssGSEA score vector and spatial coordinate
to a 256-d embedding $\mathbf{h}_\text{path}$.
UNI's 1024-d CLS token is projected by a learned linear head to the same space:
\begin{equation}
    \mathbf{h}_\text{img} = W_\text{proj}\,\mathrm{UNI}(P_s) \;\in\; \mathbb{R}^{256}.
\end{equation}
Both embeddings are L2-normalized and aligned via symmetric NT-Xent loss with a
learnable log-parameterized temperature $\tau$ clamped to $[10^{-2}, 1]$:
\begin{equation}
    \mathcal{L}_\text{CL} = -\frac{1}{2N}\sum_{i=1}^{N}
    \!\Bigl[
    \log\tfrac{e^{\mathbf{h}_i^\top \mathbf{h}_i'/\tau}}
              {\sum_j e^{\mathbf{h}_i^\top \mathbf{h}_j'/\tau}}
    + \log\tfrac{e^{\mathbf{h}_i'^\top \mathbf{h}_i/\tau}}
               {\sum_j e^{\mathbf{h}_i'^\top \mathbf{h}_j/\tau}}
    \Bigr].
\end{equation}

\textbf{Stage~2 (Supervised Prediction).}
The pathway encoder is frozen.
Two heads $f_\text{path}$, $f_\text{gene}$ are trained on the frozen image
embeddings:
\begin{equation}
    \mathcal{L}_\text{sup} =
    \|f_\text{path}(\mathbf{h}_\text{img}) - \mathbf{y}_\text{path}\|_2^2
    +
    \|f_\text{gene}(\mathbf{h}_\text{img})  - \mathbf{y}_\text{gene}\|_2^2.
\end{equation}

\subsection{Encoder Architectures}

\textbf{Pathway Encoder.}
A two-stream architecture: spatial coordinates $\mathbf{c} \in [0,1]^2$
pass through a two-layer MLP ($2\!\to\!256\!\to\!512$); pathway scores pass
through a linear projection ($P\!\to\!512$).
The sum is processed by a 2-layer Transformer encoder (hidden dim 512, 8 heads,
feedforward dim 1024) and projected to $\mathbb{R}^{256}$.

\textbf{Vision Encoder.}
UNI (DINOv2 ViT-Large/16, trained on 100{,}000 H\&E WSIs~\cite{uni2024})
produces 1024-d CLS embeddings projected to 256-d.
During Stage~1, the last four transformer blocks of UNI are unfrozen and
updated jointly with all other trainable parameters; the backbone prefix is kept
frozen throughout training.

\subsection{MLP Head (Baseline)}
\label{ssec:mlp}

The baseline follows PEaRL~\cite{pearl2025}: a two-layer ReLU MLP
\begin{equation}
    f_\text{MLP}(\mathbf{h}) =
    W_2\,\mathrm{ReLU}(W_1\mathbf{h} + \mathbf{b}_1) + \mathbf{b}_2,
\end{equation}
where $W_1 \in \mathbb{R}^{256\times 256}$ and $W_2 \in \mathbb{R}^{D_\text{out}\times 256}$.
Trained end-to-end via AdamW on $\mathcal{L}_\text{sup}$.

\subsection{TabPFN Head (Proposed)}
\label{ssec:tabpfn}

Figure~\ref{fig:tabpfn_arch} shows our proposed framework.
We replace both MLP heads with a bank of \texttt{TabPFNRegressor} instances
operating in \textit{pure} mode: one independent regressor per output dimension,
covering \emph{all} $D_\text{out}$ gene and pathway dimensions.

\begin{figure}[t]
\centering
\includegraphics[width=\columnwidth]{Screenshot 2026-05-14 at 11.03.51 PM.png}
\caption{Proposed PEaRL+TabPFN framework.
Identical to the original PEaRL (Fig.~\ref{fig:pearl_arch}) except that the
Gene MLP and Pathway MLP heads at inference are replaced by Gene TabPFN and
Pathway TabPFN regressor banks respectively.
The shared encoders and contrastive pretraining are unchanged.}
\label{fig:tabpfn_arch}
\end{figure}

\textbf{Fitting procedure:}
After Stage~2 completes (training the MLP heads to finalize the image-side
projection head), we extract training image embeddings
$\mathbf{H}_\text{train} \in \mathbb{R}^{N_\text{train}\times 256}$
via a single frozen-encoder forward pass.
For each dimension $d$, a \texttt{TabPFNRegressor} (32 estimators) is fitted
on $(\mathbf{H}_\text{train},\, \mathbf{y}_{:,d}^\text{train})$ ---
an in-context forward pass requiring no gradient computation.

\textbf{Inference.}
Each validation embedding is passed through all $D_\text{out}$ fitted regressors
in sequence.
Unlike the MLP's shared hidden layer, each TabPFN regressor is fully independent;
cross-dimension dependencies can only be recovered implicitly through the shared
256-d input embedding.

% ============================================================================
\section{Experimental Setup}
\label{sec:setup}
% ============================================================================

\subsection{Datasets}

We evaluate on two HEST-1k~\cite{hest1k} cohorts (Table~\ref{tab:datasets}).
The lymph node cohort (74{,}220 spots) exceeds TabPFN's practical in-context
row limit of $\sim$10{,}000 and is excluded from this comparison.

\begin{table}[h]
\centering
\caption{HEST-1k evaluation cohorts. Lymph is excluded due to TabPFN's
row-count limit ($\sim$10k spots).}
\label{tab:datasets}
\begin{tabular}{lccc}
\toprule
\textbf{Cohort} & \textbf{Sections} & \textbf{Spots} & \textbf{Pathways} \\
\midrule
Breast & 36 & 13{,}620 & 775 \\
Skin   & 12 &  8{,}671 & 609 \\
Lymph  & 24 & 74{,}220 & 1{,}100 \\
\bottomrule
\end{tabular}
\end{table}

\subsection{Preprocessing}

\textbf{Genes.}
Genes detected in fewer than 1{,}000 spots are removed before normalization.
Counts are CP10K-normalized, log-$(1{+}x)$-transformed, and subjected to
8-nearest-neighbor spatial smoothing.
The top 1{,}000 HVGs are selected by Scanpy's Seurat-flavor dispersion method
and per-gene min-max scaled to $[0,1]$.

\textbf{Pathways.}
ssGSEA scores are computed against Reactome (human pathways, 5--500 genes) and
MSigDB Hallmark (50 curated oncology pathways).
Scores from all sections are pooled and the top-$P$ pathways by pooled
standard deviation are retained.
Pathway targets are used at raw ssGSEA scale (not z-normalized).

\textbf{Images.}
$224\!\times\!224$ RGB patches are ImageNet-normalized before the UNI forward pass.

\subsection{Training and Evaluation Protocol}

Figure~\ref{fig:exp_config} summarizes the experimental configuration.
Both heads share all training settings: AdamW (lr $= 10^{-4}$,
weight decay $= 10^{-3}$, cosine annealing), up to 100 epochs per stage,
early stopping with patience 15 on validation loss.
Batch size is 8; max spots per section is 400 (24~GB GPU memory constraint).
All experiments use 5-fold cross-validation with an 80/20 spot-level
train-validation split and seed 42.

\begin{figure}[h]
\centering
\includegraphics[width=\columnwidth]{fig_exp_config.png}
\caption{Experiment configuration. \textbf{Left (TabPFN)}: TabPFN v2, single
in-context fitting pass with 32 estimators per dimension, no gradient updates.
\textbf{Right (MLP)}: 2-layer feedforward network trained with MSE loss,
AdamW optimizer, 100 epochs. \textbf{Center (Shared)}: all other settings are
identical across both heads.}
\label{fig:exp_config}
\end{figure}

Evaluation metrics are PCC (Pearson correlation on the flattened
$N_\text{val}\!\times\!D$ tensor, constant columns filtered), MSE, and MAE.

% ============================================================================
\section{Results}
\label{sec:results}
% ============================================================================

\subsection{Breast Cancer}

Figure~\ref{fig:results_breast} shows the quantitative results on the Breast
cohort.
Table~\ref{tab:breast} reproduces the numbers in LaTeX for reference.

\begin{figure}[t]
\centering
\includegraphics[width=\columnwidth]{Screenshot 2026-05-14 at 10.54.08 PM.png}
\caption{Results on HEST-1k Breast (36 sections, 13{,}620 spots, 5-fold CV,
80/20 spot-level split, $\leq$400 spots per section).
Best result per metric in each panel is in bold.
``PEaRL (original)'' = published numbers from~\cite{pearl2025}.
``PEaRL (reproduced)'' = our re-implementation under the stated conditions.}
\label{fig:results_breast}
\end{figure}

\begin{table}[h]
\centering
\caption{Breast cohort quantitative summary.}
\label{tab:breast}
\setlength{\tabcolsep}{4pt}
\begin{tabular}{lccc ccc}
\toprule
& \multicolumn{3}{c}{\textbf{Gene}} & \multicolumn{3}{c}{\textbf{Pathway}} \\
\cmidrule(lr){2-4}\cmidrule(lr){5-7}
\textbf{Method} & PCC$\uparrow$ & MSE$\downarrow$ & MAE$\downarrow$
                & PCC$\uparrow$ & MSE$\downarrow$ & MAE$\downarrow$ \\
\midrule
PEaRL (orig.)~\cite{pearl2025}
  & 0.5868 & 0.0732 & 0.1828 & 0.5055 & 0.0017 & 0.0314 \\
PEaRL+MLP (ours)
  & 0.5196 & 0.0139 & 0.0876 & \textbf{0.6625} & 0.0215 & 0.1240 \\
PEaRL+TabPFN
  & \textbf{0.5203} & \textbf{0.0128} & \textbf{0.0820}
  & 0.6607 & \textbf{0.0185} & \textbf{0.1050} \\
\bottomrule
\end{tabular}
\end{table}

TabPFN achieves higher gene PCC ($+0.14\%$) and lower gene MSE ($-8.0\%$,
$0.0128$ vs.\ $0.0139$) and MAE ($-6.4\%$) than the reproduced MLP.
On pathway prediction, the MLP attains marginally higher PCC ($0.6625$
vs.\ $0.6607$), while TabPFN reduces pathway MSE by 13.9\% and MAE by 15.3\%,
reflecting more calibrated magnitude predictions.

\subsection{Skin Cancer}

Figure~\ref{fig:results_skin} shows results on the Skin cohort.

\begin{figure}[t]
\centering
\includegraphics[width=\columnwidth]{Screenshot 2026-05-14 at 10.54.24 PM.png}
\caption{Results on HEST-1k Skin (12 sections, 8{,}671 spots, 5-fold CV).
Conditions identical to Breast except dataset size (609 pathways).}
\label{fig:results_skin}
\end{figure}

\begin{table}[h]
\centering
\caption{Skin cohort quantitative summary.}
\label{tab:skin}
\setlength{\tabcolsep}{4pt}
\begin{tabular}{lccc ccc}
\toprule
& \multicolumn{3}{c}{\textbf{Gene}} & \multicolumn{3}{c}{\textbf{Pathway}} \\
\cmidrule(lr){2-4}\cmidrule(lr){5-7}
\textbf{Method} & PCC$\uparrow$ & MSE$\downarrow$ & MAE$\downarrow$
                & PCC$\uparrow$ & MSE$\downarrow$ & MAE$\downarrow$ \\
\midrule
PEaRL (orig.)~\cite{pearl2025}
  & 0.3756 & 0.1405 & 0.2726 & 0.3523 & 0.0029 & 0.0404 \\
PEaRL+MLP (ours)
  & 0.8691 & 0.1074 & 0.1841 & \textbf{0.4982} & \textbf{0.8163} & 0.7281 \\
PEaRL+TabPFN
  & \textbf{0.8704} & \textbf{0.1065} & \textbf{0.1831}
  & 0.4920 & 0.8252 & \textbf{0.7230} \\
\bottomrule
\end{tabular}
\end{table}

On Skin, TabPFN achieves gene PCC of $0.8704$ vs.\ $0.8691$ ($+0.15\%$)
with a commensurate reduction in gene MSE ($-0.9\%$).
The substantially higher PCCs relative to the original paper's $0.3756$ reflect
the easier spot-level split task vs.\ the paper's section-stratified evaluation.

\subsection{Per-Dimension PCC Analysis}

Global flatten PCC is dominated by the high-variance output dimensions.
Figure~\ref{fig:per_dim_pcc} shows the full per-dimension PCC distributions
for both heads across all gene (1{,}000 dims) and pathway (775 dims) targets.

\begin{figure}[h]
\centering
\includegraphics[width=\columnwidth]{fig_per_dim_pcc.png}
\caption{Per-dimension PCC distributions on the Breast cohort.
Dashed lines mark distribution means.
TabPFN narrows the left tail of the gene distribution (fewer negative-correlation
dimensions), reflecting its regularized in-context prior suppressing overfitting
on low-variance targets.
The MLP maintains a marginal pathway PCC advantage (higher mean) but a broader
spread.}
\label{fig:per_dim_pcc}
\end{figure}

TabPFN produces a narrower left tail in the gene PCC distribution, indicating
more consistent predictions across low-variance targets where the MLP tends to
overfit.
For pathways, the MLP's higher mean PCC ($0.6625$ vs.\ $0.6607$) is accompanied
by a broader distribution --- the MLP's shared hidden layer captures inter-pathway
correlations while TabPFN's per-dimension independence cannot.


% ============================================================================
\section{Discussion}
\label{sec:discussion}
% ============================================================================

\subsection{Why TabPFN Transfers to Pathology Embeddings}

PEaRL's Stage-1 contrastive pretraining imposes a specific geometric structure
on image embeddings: L2-normalized 256-d vectors clustered by pathway
co-expression similarity.
This compact, norm-constrained representation reduces effective input
dimensionality relative to raw pixel space, bringing the head task into
TabPFN's operating regime.
At $\sim$320 training spots per fold (80\% of 400), the sample size is also
well within TabPFN's published sweet spot ($N \leq 10{,}000$), where its
meta-learned prior provides the strongest regularization over gradient-trained
networks.

TabPFN achieves competitive or better performance with \emph{zero training epochs}
vs.\ 100-epoch MLP across both tissue types.
This demonstrates that PEaRL's 256-d embeddings form structured tabular data
naturally compatible with TabPFN's Bayesian inductive biases, which transfer
effectively to small biological datasets.

\subsection{Gene vs.\ Pathway Prediction Behavior}

TabPFN's advantage is stronger at the gene level than the pathway level.
Two factors contribute:
(1)~\textit{Target independence}: gene expression dimensions are less
biologically correlated than pathway scores; the MLP's shared layer provides a
stronger inductive bias for capturing inter-pathway dependencies that TabPFN's
per-dimension independence cannot exploit.
(2)~\textit{Normalization sensitivity}: pathway MSE comparisons are sensitive
to target scaling. Under z-normalized pathways (unit variance), absolute MSE
values are $\approx$400$\times$ larger than under raw ssGSEA scaling, shifting
which method appears better on absolute-error metrics while leaving PCC
scale-invariant.

\subsection{Compute Trade-offs}

Table~\ref{tab:compute} summarizes the compute budget.
The TabPFN fitting phase (materializing 1{,}775 independent regressors per fold)
costs approximately $15\times$ the MLP Stage-2 training time.
At inference, a single batched predict call per dimension makes TabPFN only
marginally slower than a single MLP forward pass.

\begin{table}[h]
\centering
\caption{Approximate compute on a 24~GB GPU, Breast cohort, 5-fold CV.}
\label{tab:compute}
\begin{tabular}{lcc}
\toprule
\textbf{Phase} & \textbf{MLP Head} & \textbf{TabPFN Head} \\
\midrule
Stages 1+2 (shared)        & $\sim$7 hr   & $\sim$7 hr \\
TabPFN fitting              & ---           & $\sim$45 hr \\
\midrule
\textbf{Total}              & $\sim$7 hr   & $\sim$52 hr \\
\bottomrule
\end{tabular}
\end{table}

The $\sim$7$\times$ overhead is justified when zero-hyperparameter head fitting
and built-in uncertainty quantification (TabPFN's Bayesian posterior) are
operationally valuable.
In throughput-sensitive production settings, the trained MLP remains more
practical.

\subsection{Reproduced vs.\ Original PEaRL Numbers}

The reproduced MLP achieves pathway PCC of $0.66$ vs.\ the original paper's
$0.51$ on Breast, and gene PCC of $0.52$ vs.\ $0.59$.
The pathway PCC gap is primarily attributable to the spot-level 80/20 split,
which allows within-section train-validation overlap (an easier task) vs.\ the
paper's section-stratified split.
The gene PCC gap ($0.52$ vs.\ $0.59$) likely reflects the paper's use of the
gated UNI checkpoint with full last-4-block fine-tuning vs.\ our ImageNet ViT
baseline, and additional preprocessing differences.

% ============================================================================
\section{Limitations}
\label{sec:limitations}
% ============================================================================

\textbf{Spot-count cap.}
All runs use $\leq$400 spots per section (memory constraint); a full-cohort
run with more spots would increase effective training set size and potentially
favor the MLP at the boundary of TabPFN's row limit.

\textbf{Spot-level split.}
The 80/20 spot-level split allows within-section train-val leakage.
A section-stratified split (our planned protocol)
would evaluate stricter cross-patient generalization.

\textbf{Per-dimension independence.}
TabPFN cannot model joint target distributions.
This limits co-expression structure recovery, which may matter for downstream
pathway clustering and survival analysis.

\textbf{Lymph node cohort.}
With 74{,}220 spots, Lymph exceeds TabPFN's practical row limit and requires
heavy subsampling that would confound a fair comparison.

\textbf{UNI backbone access.}
These results use the ImageNet ViT-Large backbone (timm) due to HuggingFace
gating constraints in the compute environment.
Runs with the true UNI checkpoint may yield different absolute PCCs;
the relative ordering between heads is expected to be stable.

% ============================================================================
\section{Future Work}
\label{sec:future}
% ============================================================================

\textbf{Uncertainty quantification.}
TabPFN provides per-spot predictive uncertainties without additional computation,
enabling spatial maps of prediction confidence --- a capability not available
from a deterministic MLP without MC dropout or ensembling.

\textbf{Ensemble head.}
A convex combination of MLP and TabPFN predictions could capture both the MLP's
cross-dimension correlation structure and TabPFN's per-dimension calibration.

\textbf{Section-stratified full evaluation.}
The planned  protocol (5-fold GroupKFold by
section, gated UNI backbone, raw ssGSEA targets) will enable direct comparison
against the original PEaRL paper numbers.

\textbf{Alternative tabular methods.}
Extending the comparison to XGBoost, CatBoost, and SAINT~\cite{saint2021}
would situate TabPFN within the broader tabular regression landscape.

\textbf{Pan-cancer generalization.}
Evaluating on lung, colorectal, and prostate HEST-1k cohorts would test
generalization across tissue architectures.

% ============================================================================
\section{Conclusion}
\label{sec:conclusion}
% ============================================================================

We presented a controlled evaluation of \textsc{TabPFN} as a drop-in
prediction head in the PEaRL spatial transcriptomics framework, replacing
gradient-trained MLP heads with a bank of pretrained in-context tabular
regressors on both Breast and Skin cancer cohorts from HEST-1k.
TabPFN matches or exceeds the MLP on gene PCC and consistently reduces MSE
and MAE --- with \emph{zero gradient updates} at the head level.
Per-dimension PCC analysis reveals that TabPFN's gains arise from improved
per-dimension calibration, narrowing the left tail of low-PCC dimensions
where the MLP tends to overfit on low-variance targets.
These results demonstrate that PEaRL's contrastive pretraining produces
256-dimensional embeddings that are structurally compatible with TabPFN's
Bayesian inductive prior, and that the in-context learning paradigm is a
viable, training-free alternative for small-sample spatial transcriptomics
prediction.

% ============================================================================
\section*{Acknowledgment}
% ============================================================================

This work was completed as part of ME~5920 (Data Analytics and Machine Learning for Cyber-Physical Systems Applications) under the supervision of Dr. Soumik Sarkar. The authors also thank Koushik Howlader for valuable discussions and insights regarding the pathway design and methodological direction of this work.


% ============================================================================
\bibliographystyle{IEEEtran}
\begin{thebibliography}{99}

\bibitem{pearl2025}
S.~Majumder, S.~Kapse, M.~Bhattacharya, X.~Xu, A.~Yurovsky,
and P.~Prasanna,
``PEaRL: Pathway-Enhanced Representation Learning for Gene and Pathway
Expression Prediction from Histology,''
\textit{Proc.\ IEEE/CVF Winter Conf.\ on Applications of Computer Vision
(WACV)}, 2026.
[\href{https://arxiv.org/abs/2510.03455}{arXiv:2510.03455}]

\bibitem{hest1k}
G.~Jaume, P.~Doucet, A.~H.~Song, M.~Y.~Lu, \textit{et~al.},
``HEST-1k: A Dataset for Spatial Transcriptomics and Histology Image
Analysis,''
\textit{Advances in Neural Information Processing Systems (NeurIPS)},
vol.~37, pp.~53798--53833, 2024.
[\href{https://arxiv.org/abs/2406.16192}{arXiv:2406.16192}]

\bibitem{uni2024}
R.~J.~Chen, T.~Ding, M.~Y.~Lu, \textit{et~al.},
``Towards a General-Purpose Foundation Model for Computational Pathology,''
\textit{Nature Medicine}, vol.~30, no.~3, pp.~850--862, 2024.
[\href{https://doi.org/10.1038/s41591-024-02857-3}{doi:10.1038/s41591-024-02857-3}]

\bibitem{tabpfn2023}
K.~Eggensperger, M.~Feurer, N.~Hollmann, S.~M\"{u}ller,
and F.~Hutter,
``Accurate Predictions on Small Data with a Tabular Foundation Model,''
\textit{Nature}, vol.~637, no.~8045, pp.~319--326, 2025.
[\href{https://doi.org/10.1038/s41586-024-08328-6}{doi:10.1038/s41586-024-08328-6}]

\bibitem{conch2024}
M.~Y.~Lu, B.~Chen, D.~F.~K.~Williamson, R.~J.~Chen, \textit{et~al.},
``A Visual-Language Foundation Model for Computational Pathology,''
\textit{Nature Medicine}, vol.~30, no.~3, pp.~863--874, 2024.
[\href{https://doi.org/10.1038/s41591-024-02856-4}{doi:10.1038/s41591-024-02856-4}]

\bibitem{ssgsea}
D.~A.~Barbie, P.~Tamayo, J.~S.~Boehm, S.~Y.~Kim, \textit{et~al.},
``Systematic RNA Interference Reveals That Oncogenic KRAS-Driven Cancers
Require TBK1,''
\textit{Nature}, vol.~462, pp.~108--112, 2009.
[\href{https://doi.org/10.1038/nature08460}{doi:10.1038/nature08460}]

\bibitem{stnet}
B.~He, L.~Bergenstrahle, L.~Stenbeck, \textit{et~al.},
``Integrating Spatial Gene Expression and Breast Tumour Morphology via
Deep Learning,''
\textit{Nature Biomedical Engineering}, vol.~4, no.~8, pp.~827--834, 2020.
[\href{https://doi.org/10.1038/s41551-020-0578-x}{doi:10.1038/s41551-020-0578-x}]

\bibitem{his2st}
Y.~Zeng, Z.~Wei, W.~Yu, R.~Yin, Y.~Yuan, B.~Li, Q.~Huang,
Y.~Wang, and G.~Gao,
``Spatial Transcriptomics Prediction from Histology Jointly Through
Transformer and Graph Neural Networks,''
\textit{Briefings in Bioinformatics}, vol.~23, no.~5, 2022.
[\href{https://doi.org/10.1093/bib/bbac297}{doi:10.1093/bib/bbac297}]

\bibitem{bleep}
R.~Xie, R.~Zheng, J.~Wang, \textit{et~al.},
``Spatially Resolved Gene Expression Prediction from Histology Images
via Bimodal Contrastive Learning,''
\textit{Advances in Neural Information Processing Systems (NeurIPS)},
vol.~36, pp.~70626--70637, 2023.
[\href{https://arxiv.org/abs/2306.01859}{arXiv:2306.01859}]

\bibitem{mclstexp}
W.~Min, J.~Shi, J.~Zhang, K.~Wan, and S.~Wang,
``Multimodal Contrastive Learning for Spatial Gene Expression Prediction
Using Histology Images,''
\textit{Briefings in Bioinformatics}, vol.~25, no.~6, 2024.
[\href{https://doi.org/10.1093/bib/bbae578}{doi:10.1093/bib/bbae578}]

\bibitem{saint2021}
G.~Somepalli, M.~Goldblum, A.~Schwarzschild, C.~B.~Bruss,
and T.~Goldstein,
``SAINT: Improved Neural Networks for Tabular Data via Row Attention
and Contrastive Pre-Training,''
\textit{arXiv preprint}, 2021.
[\href{https://arxiv.org/abs/2106.01342}{arXiv:2106.01342}]



\end{thebibliography}

\end{document}
